\subsection{Kritische Würdigung}
\label{subsec:kritische-wuerdigung}

Die projektnahe Fallstudie ist begrenzt generalisierbar. Dokumentation,
Abnahme und lokale Prüfung stammen aus dem Projektkontext; die
Entwicklerperspektive kann die Evidenzauswahl beeinflussen. Literatur und
offizielle Dokumentation ersetzen keine unabhängige Replikation.
Commitgebundene Befunde bleiben dadurch gültig, ihre Übertragung auf andere
Stände oder Projekte aber begrenzt.

\paragraph{Konstruktvalidität.}
Code-LOC, Funktions- und Klassengröße sowie Komplexität bilden ausgewählte
Strukturmerkmale ab. Sie messen weder fachliche Verständlichkeit, Kohäsion noch
korrekte Verantwortungsverteilung. Importregeln erfassen nicht jede
dynamische, transitive oder semantische Architekturverletzung. Insbesondere
900 Code-LOC sind eine großzügige projektspezifische Obergrenze und kein
universelles Wartbarkeitskriterium.

\paragraph{Interne Validität.}
Die Auswertung stammt aus dem Projekt selbst. Fixtures für Grenzwerte und
Architektur können auf die implementierten Regeln zugeschnitten sein.
Nachgebildete Ruff- und ESLint-Ausgaben sowie die fehlende
negative Rust-Fixture begrenzen den Integrationsnachweis. Hinzu kommen fehlende
integrierte Tests für die 901-LOC-Ausnahme bis zum Gate-Status sowie für
absolute und ausgeschlossene Exception-Pfade. Bestandene Tests schließen
Fehler außerhalb der geprüften Fälle nicht aus.

\paragraph{Externe Validität.}
Untersucht wird ein einzelnes profilbasiertes Template. Grenzwerte und
Architekturabbildung sind nicht ohne Weiteres auf kleine Starter, große
Enterprise-Systeme oder generierte Codebasen übertragbar. Übertragbar sind
eher die Prinzipien zentraler Policy, deklarativer Variabilität und gemeinsamer
Prüfeinstiege als die konkreten Zahlenwerte.

\paragraph{Reliabilität.}
Ergebnisse hängen von Commit, Toolversionen, Betriebssystem, Paketregistries
und nativen Bibliotheken ab. Die abweichenden Master- und Profilergebnisse
sowie der lokal fehlgeschlagene Rust-/Tauri-Pfad zeigen diese Abhängigkeit.
Ohne Coverage-Messung, Compose-Start und erfolgreiche Plattformläufe entsteht
kein umfassender Qualitäts- oder Produktionsnachweis.

\paragraph{Langzeitwirkung.}
Die Untersuchung belegt die technische Erkennung definierter Verstöße zum
Beobachtungszeitpunkt. Wartungsaufwand, Defektrate, Refactoring-Häufigkeit,
Entwicklungszeit und Modellbefolgung wurden nicht longitudinal gemessen. Das
Quality Gate garantiert daher weder Clean Code noch eine kausale Verbesserung
langfristiger Wartbarkeit.
